% A_data_predictors.tex -- online appendix: construction of the data and the predictors (2026-10-26, integration pass
% for JFEc; DATA agent). Material moved out of Sections 3 and 4 (rehaul/03_data.tex, rehaul/04_features.tex), lightly
% edited, with the audit fixes of 2026-10-25 and the author decisions D13 (FOMC fixes) and D15 (holiday sessions) of
% writeup/INTEGRATION_PLAN_2026-10-26.md. One \subsection per block, each with an app:rh_dp_* label; the APPENDIX agent
% wires this file in (suggested: inside \section{Data and predictors}, label app:rh_data, of rehaul/10_appendix.tex).
% House rules: every number from a generated macro; replaced text parked in dated comments in the section files.
%
% Stop-gap definitions (2026-10-26, DATA agent): constants of the code with no generated macro yet, at the values of the
% code. \providecommand, so the generators' macros win once they exist (requested from the NUMBERS agent; generator in
% brackets). The first two and the FOMC dates repeat the lines of 04_features.tex and 03_data.tex (no-ops here then).
% 2026-10-26 (coordinator): parked stop-gap, the macro is now generated -- \providecommand{\rhdmWinsorWindow}{240}  % [make_rehaul_data_methods.py] target.py rolling_winsorize window (robust_transform default 240)
% 2026-10-26 (coordinator): parked stop-gap, the macro is now generated -- \providecommand{\rhdmIqrFloorFrac}{0.01}  % [make_rehaul_data_methods.py] scaling.py IQR_FLOOR_FRAC
% 2026-10-26 (coordinator): parked stop-gap, the macro is now generated -- \providecommand{\rhdmWinsorLoPct}{1}  % [make_rehaul_data_methods.py] target.py WINSOR_LOWER_Q x 100
% 2026-10-26 (coordinator): parked stop-gap, the macro is now generated -- \providecommand{\rhdmWinsorHiPct}{99}  % [make_rehaul_data_methods.py] target.py WINSOR_UPPER_Q x 100
% 2026-10-26 (coordinator): parked stop-gap, the macro is now generated -- \providecommand{\rhdmBandLo}{0.5}  % [make_rehaul_data_methods.py] target.py DIURNAL_SLOT_BAND_FRAC (band [f rho, rho / f])
% 2026-10-26 (coordinator): parked stop-gap, the macro is now generated -- \providecommand{\rhdmBandHi}{2}  % [make_rehaul_data_methods.py] 1 / DIURNAL_SLOT_BAND_FRAC
% 2026-10-26 (coordinator): parked stop-gap, the macro is now generated -- \providecommand{\rhdmSlotRefWindow}{250}  % [make_rehaul_data_methods.py] target.py DIURNAL_SLOT_REF_WINDOW
% 2026-10-26 (coordinator): parked stop-gap, the macro is now generated -- \providecommand{\rhdmSlotRefMin}{20}  % [make_rehaul_data_methods.py] target.py DIURNAL_SLOT_REF_MIN
% 2026-10-26 (coordinator): parked stop-gap, the macro is now generated -- \providecommand{\rhdmSdFloorFrac}{0.1}  % [make_rehaul_data_methods.py] target.py DIURNAL_STD_FLOOR_FRAC
% 2026-10-26 (coordinator): parked stop-gap, the macro is now generated -- \providecommand{\rhdmLogFloor}{\ensuremath{10^{-12}}}  % [make_rehaul_data_methods.py] target.py log rule, log(clip(x, 1e-12))
% 2026-10-26 (coordinator): parked stop-gap, the macro is now generated -- \providecommand{\rhdmAsinhScaleFrac}{0.01}  % [make_rehaul_data_methods.py] target.py asinh scale floor, 0.01 x expanding median
% 2026-10-26 (coordinator): parked stop-gap, the macro is now generated -- \providecommand{\rhdmAsinhHalfLife}{12{,}000}  % [make_rehaul_data_methods.py] ASINH_SCALE_HALFLIFE_DAYS x PERIODS_PER_DAY
% 2026-10-26 (coordinator): parked stop-gap, the macro is now generated -- \providecommand{\rhdmAsinhHalfLifeDays}{250}  % [make_rehaul_data_methods.py] target.py ASINH_SCALE_HALFLIFE_DAYS
% 2026-10-26 (coordinator): parked stop-gap, the macro is now generated -- \providecommand{\rhdmIqrGridStep}{480}  % [make_rehaul_data_methods.py] executor.py _expanding_real_iqr grid_step
% 2026-10-26 (coordinator): parked stop-gap, the macro is now generated -- \providecommand{\rhdmIqrMinRows}{8}  % [make_rehaul_data_methods.py] executor.py _expanding_real_iqr, rows before a value
% 2026-10-26 (coordinator): parked stop-gap, the macro is now generated -- \providecommand{\rhdmIqrEps}{\ensuremath{10^{-12}}}  % [make_rehaul_data_methods.py] scaling.py IQR guard (below it -> 1)
% 2026-10-26 (coordinator): parked stop-gap, the macro is now generated -- \providecommand{\rhdmBarsPerSession}{46}  % [make_rehaul_data_methods.py] executor.py BARS_PER_SESSION (ERODE_IN_BARS)
% 2026-10-26 (coordinator): parked stop-gap, the macro is now generated -- \providecommand{\rhdsFomcSchedMeeting}{17--18 March 2020}  % [make_rehaul_disclosures.py] loading.py FOMC_2020_SCHEDULED
% 2026-10-26 (coordinator): parked stop-gap, the macro is now generated -- \providecommand{\rhdsFomcReplaced}{15 March 2020}  % [make_rehaul_disclosures.py] loading.py FOMC_2020_ANNOUNCEMENT
% 2026-10-26 (coordinator): parked stop-gap, the macro is now generated -- \providecommand{\rhdsFomcNextIssued}{29 April 2020}  % [make_rehaul_disclosures.py] loading.py FOMC_2020_NEXT
% 2026-10-26 (coordinator): parked stop-gap, the macro is now generated -- \providecommand{\rhPending}{\textit{pending}}

% ----------------------------------------------------------------------------------------------------------------------
% From 03_data.tex (2026-10-26): the liquidity, market, sentiment, implied-volatility and volatility-demand bullets,
% the time stamps and the ends of the feeds. First-print years: \rhlv...First (writeup/make_rehaul_live_windows.py,
% the first print of each series) in place of the typed years; the volatility-demand end: \rhdsVdLast.
\subsection{Data sources}
\label{app:rh_dp_sources}
\begin{itemize}
% From 03_data.tex (2026-10-26): the display of the futures' realized moments and the sentence that names them.
% Sums read off the column names: the builder of data/core_stats.parquet is not in the repo
% (live/close_signal/features.py:1-25); sumpret2 = positive semivariance is an assumed reading; bipower unnormalized.
\item \textbf{Realized moments.} With $r_{t,i,k}$ the one-minute log-price changes of Section~\ref{sec:rh_data}, the
  \bkSizeMom{} realized-moment series are the one-minute sums
  % 2026-10-19 (user): raw parquet names for now (\ifrawseries, paper_rehaul.tex preamble): each sum is named by its
  % column of data/core_stats.parquet. The \else branch is the display of the symbols.
  \ifrawseries
  \begin{gather*}
    \sername{sumret}_{t,i} = \sum_k r_{t,i,k}, \quad \sername{sumabsret}_{t,i} = \sum_k |r_{t,i,k}|, \quad
    \sername{sumret3}_{t,i} = \sum_k r_{t,i,k}^3,\\
    \sername{sumret4}_{t,i} = \sum_k r_{t,i,k}^4, \quad
    \sername{sumpret2}_{t,i} = \sum_k r_{t,i,k}^2\,\mathbf{1}\{r_{t,i,k} > 0\},\\
    \sername{sumbipow}_{t,i} = \sum_{k \ge 2} |r_{t,i,k}|\,|r_{t,i,k-1}|, \quad
    \sername{sumautocov}_{t,i} = \sum_{k \ge 2} r_{t,i,k}\, r_{t,i,k-1}.
  \end{gather*}
  \else
  \[
    \sum_k r_{t,i,k}, \quad \sum_k |r_{t,i,k}|, \quad \sum_k r_{t,i,k}^3, \quad \sum_k r_{t,i,k}^4, \quad
    \mathrm{RS}^{+}_{t,i} = \sum_k r_{t,i,k}^2\,\mathbf{1}\{r_{t,i,k} > 0\},
  \]
  \[
    \mathrm{BV}_{t,i} = \sum_{k \ge 2} |r_{t,i,k}|\,|r_{t,i,k-1}|, \qquad
    \mathrm{AC}_{t,i} = \sum_{k \ge 2} r_{t,i,k}\, r_{t,i,k-1}.
  \]
  \fi
  The third sum is the unnormalized realized skewness of \citet{AmayaChristoffersenJacobsVasquez2015} (before the
  division by $\RV^{3/2}_{t,i}/\sqrt{N_{t,i}}$), the fourth the realized quarticity without the factor $N_{t,i}/3$,
  $\ser{sumpret2}{\mathrm{RS}^{+}}_{t,i}$ the positive realized semivariance of
  \citet{BarndorffNielsenKinnebrockShephard2010}, $\ser{sumbipow}{\mathrm{BV}}_{t,i}$ the bipower variation of
  \citet{BarndorffNielsenShephard2004} without its $\pi/2$ factor, and $\ser{sumautocov}{\mathrm{AC}}_{t,i}$ the
  first-order realized autocovariance. The realized variance itself is the target and not a predictor.
\item \textbf{Liquidity.} From the front-month futures contract we take the total contract volume traded in the half
  hour and the number of observed one-minute intervals $N_{t,i}$\ifrawseries{} (\sername{numobs})\fi. From intraday
  data on the S\&P~500 constituents, 1993 to 2024, we form equal- and value-weighted cross-sectional aggregates of
  share turnover, buyer-initiated turnover, seller-initiated turnover, the effective spread, the quoted spread and
  order-flow imbalance. Constituent turnover is the number of shares traded in the half hour divided by the prior
  day's shares outstanding; buyer- and seller-initiated turnover aggregate the trades by their estimated direction.
  The effective spread is the average percentage deviation of the transaction price from the prevailing mid-quote,
  and the quoted spread the average percentage bid--ask spread. Equal-weighted aggregates are simple cross-sectional
  averages; value-weighted aggregates weight each constituent by its prior-day market capitalization. For each
  weighting ($\ser{ofi_ewstock}{\mathrm{OFI}^{\mathrm{EW}}}$ and $\ser{ofi_vwstock}{\mathrm{OFI}^{\mathrm{VW}}}$),
  order-flow imbalance is
  \[
    \mathrm{OFI}_{t,i} = \frac{\mathrm{Buy}_{t,i} - \mathrm{Sell}_{t,i}}{\mathrm{Buy}_{t,i} + \mathrm{Sell}_{t,i}}
      \in [-1, 1],
    \qquad \text{missing if } \mathrm{Buy}_{t,i} + \mathrm{Sell}_{t,i} = 0 ,
  \]
  with $\mathrm{Buy}_{t,i}$ and $\mathrm{Sell}_{t,i}$ the aggregate buyer- and seller-initiated turnover of interval
  $(t,i)$.
  % loading.py:161-180, built on the unfilled legs.
\item \textbf{Market moments.} For each weighting, the \bkSizeMew{} constituent moments are the realized variance,
  realized skewness, realized quarticity, absolute variation, bipower variation and positive semivariance, computed
  from 5-minute log returns rather than the one-minute changes of the futures moments. The constituent groups omit
  the signed return sum and the realized autocovariance; the futures group omits the realized variance, which is the
  target.
\item \textbf{Sentiment.} StockTwits message streams, from \rhlvStAttFirst{} for all posts and \rhlvStSentFirst{} for
  posts tagged by their users: attention, the number of posts in the half hour; the sentiment score, $1$ when every
  tagged post is bullish and $-1$ when every tagged post is bearish; and the sentiment count, the number of posts
  tagged bullish or bearish.
  % 2026-10-26 (integration): parked -- the visible note E30, still open (the score between its endpoints); reported to
  % FRONT for rehaul/99_author_notes.tex. Was:
  % \editc{orange}{E30}{Decision: open, for the professor.\\Problem: the text gives only the $\pm 1$ endpoints of the sentiment score.\\
  % Question for the professor: what is the formula of the score between the endpoints, and what value does it take on intervals with no tagged posts? The repository receives the score precomputed (data/spy\_and\_sentiment.parquet, column stocktwits\_sentiment, first observed in June 2010); src/data/loading.py:280 reads it as a net score in $[-1, 1]$, and no script in src/ or experiments/ computes it.}
\item \textbf{Implied volatility.} The 30-day S\&P~500 implied volatility index (VIX, from \rhlvVixFirst), the
  three-month index (VIX3M, from \rhlvVixThreeFirst) and the implied volatility of the VIX (VVIX, from
  \rhlvVvixFirst), in a vendor file that ends on \rhdsVixVendorLast.
\item \textbf{Volatility demand.} From Cboe option-volume records, \rhlvVdFirst{} to \rhdsVdLast, the net vega bought by
  customers other than market makers, in \bkSizeVd{} series that cross two option universes (options on the
  S\&P~500 index and on the top S\&P~500 exchange-traded funds, and all exchange-traded options) with two sets of
  trades (all opening and closing customer trades, and opening customer trades only).
\item \textbf{Time stamps.} Every file stamps an interval by its end, in Eastern time, on a grid of half-hour stamps
  from 00:00 to 23:30: interval $(t,i)$ holds the half hour that ends at its stamp, so the closing interval $i = c$ is
  stamped 16:00.
  % 2026-10-26 (integration): parked -- the visible \todo and the note E2, still open (vendor, price field, roll day;
  % the 17:30/18:00 stamps inside CME's maintenance halt); reported to FRONT for rehaul/99_author_notes.tex. Was:
  % "\todo{name the vendor of each file, the price inside the one-minute log-price change (trade or quote), and the day
  % whose volume selects the front-month contract.}" and the box \editc{blue}{E2}{...} of 03_data.tex (parked there).
% writeup/make_rehaul_disclosures.py (data/vix_and_voldemand.parquet, data/free_feed/*_1h_yahoo.parquet,
% data/releases.parquet; src/data/loading.py SPLICE_VIX_FREE_FEED, VIX_SPLICE_SLOTS and FOMC_CALENDAR_FIX). "The same
% value to the second decimal": |Yahoo - vendor| < 0.006 index points (experiments/data_audit_stages.py).
\item \textbf{Ends of the feeds.} Our data end on \rhdsSampleEnd, but three feeds stop earlier. The file of the implied
  volatility indices ends on \rhdsVixVendorLast. For the remaining \rhdsVixSpliceDays{} trading days we splice in the
  hourly closes of the free Yahoo Finance feed for the same three indices, each written at the end of its hour and
  carried to the next half-hour stamp, on the stamps at which the file prints each index. On the stamps both sources
  carry, from \rhdsVixOverlapStart{} to \rhdsVixVendorLast, the two print the same value to the second decimal on
  \rhdsVixAgree\% (VIX), \rhdsVixThreeMAgree\% (VIX3M) and \rhdsVvixAgree\% (VVIX) of the stamps. The FOMC release
  file ends with the statement of \rhdsFomcFileLast; we add the \rhdsFomcAdded{} scheduled statements it omits before
  the end of our data, so no FOMC value is missing (Appendix~\ref{app:rh_dp_fomc}). The volatility-demand feed ends on
  \rhdsVdLast{} and is not extended: from then on its \bkSizeVd{} series are unobserved
  (Appendix~\ref{app:rh_dp_sessions}).
\end{itemize}

% ----------------------------------------------------------------------------------------------------------------------
% 2026-10-26 (D13, author: FOMC fixes adopted): src/data/loading.py FOMC_CALENDAR_FIX / FOMC_MISSING_STATEMENTS (the
% three statements after the release file), FOMC_SCHEDULE_2020_FIX (fomc_schedule_2020: until_inv counts to the
% scheduled 2020-03-18 14:30 statement on the bars after the 2020-01-29 statement that end before the 2020-03-15 17:00
% announcement; since_inv counts from the first bar after the announcement, the Sunday 18:30 reopen, to the 2020-04-29
% statement; fomc_release and fomc_day unchanged), src/features/transforms/target.py FOMC_FLAGS_IDENTITY (fomc_day: no
% log, no clip) and FOMC_PROXIMITY_NO_WINSOR (the proximity ramps skip rolling_winsorize; the log and every later step
% unchanged). Distances: add_fomc_features on the frame before the dead-session drop, i.e. the calendar-open grid (the
% weekend break cut, every other half hour kept). results/fomcfix_2026-10-26/RULE.txt (P1, P2) and STATUS.md (gates
% G1, G3, G3b). Audit 2026-10-25 correctness-7 and fidelity-7.
\subsection{FOMC calendar}
\label{app:rh_dp_fomc}
\begin{itemize}
\item \textbf{Statements.} The scheduled statement times come from the release file, which ends with the statement of
  \rhdsFomcFileLast, and from the \rhdsFomcAdded{} scheduled statements that follow it up to the end of our data. The
  statement-day indicator equals 1 on every interval of a day with a scheduled statement. The distances
  $n^{+}_{t,i}$ and $n^{-}_{t,i}$ count half-hour intervals on the grid less the weekend break.
\item \textbf{The pre-announced calendar in March 2020.} The release file holds the statements that were issued, and
  the meeting scheduled for \rhdsFomcSchedMeeting{} is not among them: the Federal Reserve replaced it with the
  action announced on Sunday, \rhdsFomcReplaced. The proximity measures follow the calendar that was public at each
  interval. Until that announcement, $n^{+}_{t,i}$ counts the intervals to that meeting's scheduled statement; from
  the announcement on, $n^{+}_{t,i}$ counts to the next statement issued, on
  \rhdsFomcNextIssued, and $n^{-}_{t,i}$ counts from the first interval after the announcement. No statement is added,
  so the statement-day indicator is unchanged.
\item \textbf{No winsorization.} The statement-day indicator passes through every step unchanged. The proximity
  measures, after the logarithm of step~2, are deterministic ramps that rise or fall monotonically between
  statements, so each value is an extreme of its own trailing window and a rolling quantile clip would bind on most
  intervals, holding the series behind its true value. They skip the winsorization of step~3; every other step
  applies.
\item Both rules change only the two proximity series. Forecasts built before them (some robustness and ablation
  arms, marked where they are reported) count to the next statement issued in early 2020 and clip the proximity
  measures.
\end{itemize}

% ----------------------------------------------------------------------------------------------------------------------
% From 03_data.tex (2026-10-26): the early closes, the holiday sessions (D15, author 2026-10-26: keep and disclose;
% audit fidelity-4, numbers-7, correctness-5, design-11), the domain rule, the availability indicator, the fill cap and
% the no-observation counts.
% Holidays: results/fomcfix_2026-10-26/RULE.txt addendum A and holiday_sessions.csv (src/data/loading.py
% NYSE_HOLIDAY_FUTURES_SESSIONS: a date with a bar ending 10:00 .. 16:00 on the grid, no 16:00 bar, not an NYSE early
% close, and no print of any constituent series; every one an NYSE rule-calendar holiday). The forecast and scored
% counts: \rhorNTreeBars (each pooled arm) and \rhorNIntB (Table 1), whose difference is the holiday bars of the span.
% src/backtest/executor.py DROP_EARLY_CLOSE_SESSIONS, B_REGULAR_BARS_ONLY (target only: baseline_exclude =
% early_close_post_bars, bars ending 13:30 .. 17:00), DROP_HOLIDAY_SESSIONS (default off; the robustness cells C, D).
% D15, NOTE FOR OFFICIAL PUBLICATION: rerun with DROP_HOLIDAY_SESSIONS on (all models) before the final version.
\subsection{Sessions, cleaning and missing values}
\label{app:rh_dp_sessions}
\begin{itemize}
\item \textbf{Grid.} The grid $\mathcal G$ starts at the first observed realized variance; the weekend break it leaves
  out runs from Friday 20:00 to Sunday 18:30~ET:
  \[
    (t,i) \in \mathcal{G} \iff \RV_{t,i} \text{ observed and } (t,i) \notin (\text{Fri } 20{:}00,\ \text{Sun } 18{:}30) .
  \]
\item \textbf{Early closes.} The early-close sessions stay in the grid $\mathcal G$, on which every trailing average is
  computed, but the models neither estimate on nor score any interval of them, and the option-trading sample
  % 2026-10-28 (author: prune "pooled"). Was: "but the pooled models neither".
  leaves them out. The intervals after their 13:00~ET close leave the window $\mathcal D_{t,i}$ of the target's
  diurnal factor $s^2_{t,i}$; the diurnal factors of the exogenous series keep them.
\item \textbf{Holiday sessions.} On full NYSE holidays the cash market is closed and the E-mini trades a shortened
  session that ends at about 13:00~ET. Its regular-hours intervals, which carry little variance, are rows of the
  \rhorNBars-bar estimation sample and enter the windows $\mathcal D_{t,i}$ of the target's diurnal factor at their
  times of day; they are forecast but never scored, which is why each model has \rhorNTreeBars{} forecasts
  % 2026-10-28 (author: prune "pooled"). Was: "each pooled model has".
  and \rhorNIntB{} scored bars over the evaluation span. Because the estimation window counts rows, its
  \rhorWindowBars{} rows span slightly fewer than \rhorWindowDays{} trading days. In the robustness check of
  Appendix~\ref{app:rh_robust_holiday}, the holiday sessions leave the estimation sample and the windows $\mathcal D_{t,i}$, while
  every trailing average is still computed on the full grid; the linear models are refitted with their frozen
  penalties, and the every-bar comparisons of the full predictor set with the HAR benchmark change by
  \rhholEffCABarPpRng{} percentage points.
\item \textbf{Domain rule.} Before any other step, a print outside the domain of its series is set to missing: a
  nonpositive quoted or effective spread, total turnover, futures volume, observation count or index level, and a
  negative value of a series that cannot be negative. The rule removes \rhdsInvalidPrints{} prints,
  \rhdsInvalidVolume{} of them zero futures volumes on intervals that carry one-minute price changes.
\item \textbf{Availability and filling.} For each exogenous series the availability indicator $a_{t,i}$ is recorded
  before any filling. The series is then forward-filled over at most \rhdsFfillCap{} intervals, and the diurnal
  factors, the variance-stabilizing scale and the winsorization bounds are computed on the filled series; every
  interval with $a_{t,i} = 0$ then takes the imputed value of step~4 (Appendix~\ref{app:rh_dp_transform}). A gap
  longer than the cap stays missing in these statistics.
\item \textbf{No observation.} A series whose feed stops is unobserved from its last print on, and its availability
  averages fall to zero. Wherever a $k$-interval average holds no observed interval, its standardized value is set to
  zero (step~6 of Section~\ref{sec:rh_predictors}). This rule acts, on at least one interval, on
  % 2026-10-28 (author: prune "pooled"). Was: "In the pooled models this rule acts,".
  \rhdsUnobsCols{} of the \rhdsValueCols{} value averages of the full predictor set over the whole sample and on
  \rhdsUnobsColsEval{} of them within the evaluation span from \rhdsEvalStart: on every horizon of the volatility-demand
  series once their feed has stopped, and on the shorter horizons of the series that do not print on every interval
  (volatility demand, the implied volatility indices, the constituent aggregates and the StockTwits series) where the
  averaging window falls in a gap between their prints.
\end{itemize}

% ----------------------------------------------------------------------------------------------------------------------
% From 04_features.tex (2026-10-26): steps 1-4 with their displays and the activity indicator (formerly step 5), the
% closing-interval clauses moved to app:rh_dp_asfitted. Code: src/features/transforms/target.py robust_transform,
% diurnal_adjust_with_floor (CAUSAL_SD_FLOOR), series_is_signed (SIGN_BY_NAME), the map (rule list), asinh_stabilize_pair
% and the routing (CAUSAL_ASINH_ROUTING, ASINH_ROUTE_MIN_ROWS, DIURNAL_PINNED_MAX), rolling_winsorize (min_periods =
% window when allow_missing is False, as in load_and_transform's exogenous call); src/backtest/executor.py
% causal_neutral_fill (CAUSAL_NEUTRAL_FILL), causal_activity (CAUSAL_MODE_INFLATION, ACTIVE_FROM_OBSERVED,
% MODE_MIN_PRINTS, ZERO_INFLATED_FRAC), DROP_ACTIVITY_MA. Winsorization reach: audit 2026-10-25 fidelity-7 (bounds
% defined on 0% of 2014-2024 regular-hours rows for VIX3M and VVIX, 66% for VIX, 86% for the constituent aggregates;
% scratchpad audit/winsor_check.py; shares not yet macros, so not printed).
\subsection{Transformations of the exogenous series}
\label{app:rh_dp_transform}
% 2026-10-26 (coordinator: the paper keeps its bullet structure): the lead sentence as a bold-lead item.
\begin{itemize}
\item \textbf{Order.} Steps 1 to 4 of Section~\ref{sec:rh_predictors} are applied to each exogenous series $u$ in this
  order; step~5 forms the trailing averages, and step~6 (Appendix~\ref{app:rh_dp_scale}) standardizes them.
\end{itemize}
\begin{enumerate}
\item \emph{Diurnal adjustment.} Each raw series is divided by a time-of-day factor estimated over the prior-day set
  $\mathcal D_{t,i}$ of $s^2_{t,i}$ (Section~\ref{sec:rh_methods_setup}) and bounded relative to its own trailing
  median:
  \[
    \tilde u_{t,i} = \frac{u_{t,i}}{b^{(u)}_{t,i}}, \qquad
    b^{(u)}_{t,i} =
    \begin{cases}
      \clip\bigl(\operatorname{mean}_{t' \in \mathcal{D}_{t,i}} u_{t',i},\ \rhdmBandLo\,\rho_{t,i},\
        \rhdmBandHi\,\rho_{t,i}\bigr), & u \ge 0,\\[2pt]
      \max\bigl\{\clip\bigl(\operatorname{sd}_{t' \in \mathcal{D}_{t,i}} u_{t',i},\ \rhdmBandLo\,\rho_{t,i},\
        \rhdmBandHi\,\rho_{t,i}\bigr),\ \rhdmSdFloorFrac\,\bar\sigma^{(u)}_{t,i}\bigr\}, & u \text{ signed},
    \end{cases}
  \]
  \begin{itemize}
  \item where $\rho_{t,i}$ is the median of the \rhdmSlotRefWindow{} most recent unbounded factors of interval $i$
    (at least \rhdmSlotRefMin{} of them), and $\bar\sigma^{(u)}_{t,i}$ the median of the nonzero trailing standard
    deviations $\operatorname{sd}_{t'' \in \mathcal{D}_{t',i'}} u_{t'',i'}$ over all intervals $(t',i')$ up to
    $(t,i)$, every time of day pooled (before the first such standard deviation the adjusted value is undefined and
    is imputed as in step~4);
  \item a series is signed when the quantity it measures can be negative, not when its sample has a negative value
    (Table~\ref{tab:rh_dm_dictionary});
  \item $b^{(u)}_{t,i} = 1$ whenever the factor of a nonnegative series is zero or undefined; and
  \item the implied volatility indices, the StockTwits sentiment score and the FOMC series are not adjusted.
  \end{itemize}
\item \emph{Variance stabilization.} The first matching rule of
  \[
    g(v) =
    \begin{cases}
      \sqrt{v} & \text{squared returns, realized variance, turnover,}\\
        & \text{bipower variation, effective spread},\\
      \sign(v)\,|v|^{1/2} & \text{autocovariance},\\
      v^{1/3} & \text{cubed returns},\\
      \sign(v)\,|v|^{1/4} & \text{fourth-power returns},\\
      v & \text{other signed series, absolute returns},\\
      \log \max(v, \rhdmLogFloor) & \text{otherwise}
    \end{cases}
  \]
  replaces $\tilde u_{t,i}$ by $g(\tilde u_{t,i})$.
  \begin{itemize}
  \item A signed series without a power rule is routed at interval $(t,i)$ to an inverse hyperbolic sine of the raw
    series when its standard-deviation divisor sat at its floor on at least \rhdmAsinhPinnedShare\% of its intervals
    since its first print and before $(t,i)$, at least \rhdmAsinhMinRows{} of them. This routes the volatility-demand
    series, on almost all of their intervals (Table~\ref{tab:rh_dm_dictionary}), and no other series:
    \[
      \tilde u_{t,i} = \operatorname{asinh}(u_{t,i} / \kappa_{t,i}), \qquad
      \kappa_{t,i} = \max\bigl\{ e_{t,i-1},\ \rhdmAsinhScaleFrac \operatorname{med}_{j \ge 1} e_{t,i-j} \bigr\},
    \]
    where $e_{t,i}$ is an exponentially weighted moving average of $|u|$ over the nonzero intervals ($u \ne 0$) with
    a half-life of \rhdmAsinhHalfLife{} such intervals (nominally \rhdmAsinhHalfLifeDays{} days of \rhdmNIntervals{}
    intervals).
  \item The FOMC statement-day indicator passes through unchanged, and the FOMC proximity measures take the logarithm
    (Appendix~\ref{app:rh_dp_fomc}).
  \end{itemize}
\item \emph{Winsorization.} With $\mathcal U_{t,i} = \{\tilde u_{t,i-j} : j = 1, \dots, n_{\mathrm{w}}\}$ the
  preceding $n_{\mathrm{w}} = \rhdmWinsorWindow$ intervals of the grid, all times of day pooled,
  \[
    \tilde u_{t,i} \leftarrow \clip\bigl(\tilde u_{t,i},\ \underline Q^{(u)}_{t,i},\ \overline Q^{(u)}_{t,i}\bigr),
  \]
  where $\underline Q^{(u)}_{t,i}$ and $\overline Q^{(u)}_{t,i}$ are the empirical \rhdmWinsorLoPct{}st and
  \rhdmWinsorHiPct{}th percentiles of $\mathcal U_{t,i}$, the levels of the target's winsorization
  (Section~\ref{sec:rh_methods_setup}).
  \begin{itemize}
  \item The bounds are defined only when all $n_{\mathrm{w}}$ values of $\mathcal U_{t,i}$ exist, after the forward
    fill of at most \rhdsFfillCap{} intervals; otherwise the interval is not clipped. Series whose prints leave longer
    gaps are therefore clipped only where a full window of values precedes the interval: VIX3M and VVIX, which print
    only part of the day, never; the VIX and the constituent aggregates on part of the sample; the futures' own
    moments and liquidity series, which print on every interval of the grid, throughout.
  \item The FOMC series are not winsorized (Appendix~\ref{app:rh_dp_fomc}).
  \end{itemize}
\item \emph{Imputation.} An unobserved interval takes the expanding median of the series' earlier observed
  transformed values; $a_{t,i}$ enters the model separately, through its trailing averages:
  \[
    \tilde u_{t,i} \leftarrow
    \begin{cases}
      \tilde u_{t,i}, & a_{t,i} = 1,\\
      \operatorname{med}\mathcal{A}_{t,i}, & a_{t,i} = 0,\ \mathcal{A}_{t,i} \ne \emptyset,\\
      \tilde u_{(1)}, & a_{t,i} = 0,\ \mathcal{A}_{t,i} = \emptyset,
    \end{cases}
  \]
  where $\mathcal{A}_{t,i} = \{\tilde u_{t',i'} : (t',i') \prec (t,i),\ a_{t',i'} = 1\}$, $(t',i') \prec (t,i)$ means
  that $(t',i')$ is an interval of the grid strictly before $(t,i)$, and $\tilde u_{(1)}$ is the transformed value on
  the series' first observed interval. The late-start rule of step~6 keeps $\tilde u_{(1)}$ out of every forecast.
\end{enumerate}
\begin{itemize}
\item \textbf{Activity indicator (not in the pooled models).} A series is mode-inflated at interval $(t,i)$ when it has
  at least \rhdmModeMinPrints{} observed prints strictly before $(t,i)$ and their most frequent value, $m_{t,i}$,
  accounts for at least \rhdmModeShare\% of them. On every interval at which the series is mode-inflated,
  \[
    \chi_{t,i} =
    \begin{cases}
      \mathbf{1}\{u_{t,i} \ne m_{t,i}\}, & a_{t,i} = 1,\\
      \chi^{\circ}_{t,i}, & a_{t,i} = 0,
    \end{cases}
  \]
  where $\chi^{\circ}_{t,i}$ is the more frequent value of $\chi$ on the series' earlier observed intervals (0 on a
  tie), and $\chi_{t,i} = 0$ elsewhere. Under this rule \rhdmNActivePooled{} series are mode-inflated at some interval:
  the number of observations, the \bkSizeSent{} StockTwits series (on part of the sample) and the FOMC statement-day
  indicator; the volatility-demand series never are. The pooled models carry no trailing averages of $\chi$; the
  evidence requirement of the penalized linear models (Appendix~\ref{app:rh_methods_linear}) uses the same threshold of
  \rhdmModeMinPrints{} intervals, and the closing-interval panel of the earlier model set carries the averages
  (Appendix~\ref{app:rh_closing_predictors}).
\end{itemize}

% ----------------------------------------------------------------------------------------------------------------------
% From 04_features.tex (2026-10-26): step 6 in full, as the code builds every pooled design (audit 2026-10-25
% correctness-1, design-1, fidelity-1, fidelity-12): experiments/pool_linear_pretune.py run_executor(prescale=True) ->
% src/backtest/executor.py _backtest_and_save, prescale branch: _build_scale_guards (ref_iqr = _expanding_real_iqr over
% the availability mask eroded by ceil(k x rows-per-session / BARS_PER_SESSION) design rows (ERODE_IN_BARS), mode rows
% left out (CAUSAL_GUARD_MODE, _PastMode: MODE_MIN_PRINTS, ZERO_INFLATED_FRAC), grid_step 480, warm-up value 1 until 8
% rows; the *_avail_ma_* columns fixed, HAR and calendar columns ref 0), then rolling_robust_scale
% (src/features/transforms/scaling.py _rolling_scale_cols: rows [0, W) on the first window's statistics, row t >= W on
% [t - W, t); IQR floored at IQR_FLOOR_FRAC x the running maximum, < 1e-12 -> 1, then at ref_iqr), then
% late_start_warmup (LATE_START_WARMUP) and unobserved_neutral (UNOBSERVED_NEUTRAL). The masked scaler
% (apply_masked_scaling, scaling.masked_rolling_scale_col: observed rows only, every 480 rows, 512-row minimum) is called
% by no final design (only experiments/run_geometry_local.py, data_audit_stages.py, data_fix_2026_10_15.py); the design
% log of results/window_sweep_2014start_2026-10-22/w3000/run_info.json has no masked-rescaling line. The first window:
% the first W = 39,000 rows of the 13-bar design (June 1998 to May 2010), after the HAR burn-in, before every forecast.
% Live shares: writeup/make_rehaul_live_windows.py (generated/rehaul_live_numbers.tex); the training-window reading of
% audit correctness-4 stated in words (no macro yet).
\subsection{Robust standardization}
\label{app:rh_dp_scale}
\begin{itemize}
\item \textbf{Window and statistics.} The rows are those of the estimation sample, the \rhorNBars{} regular-hours bars
  of each day. Row $(t,i)$ is scaled with the median and the interquartile range of the $W = \rhorWindowBars$ rows
  before it, $\mathcal W_{t,i}$, imputed rows included, so both are recomputed at every row; the rows of the first
  window take that window's statistics, and all of them precede the first forecast. The denominator of
  Section~\ref{sec:rh_predictors}, step~6, takes the largest of three scales: the window's interquartile range; a
  relative floor, \rhdmIqrFloorFrac{} times the largest interquartile range of the column up to $(t,i)$; and the
  causal floor $\mathrm{IQR}^{\ast(k)}_{t,i}$. A scale below \rhdmIqrEps{} is replaced by 1 before the causal floor
  applies.
\item \textbf{The causal floor} $\mathrm{IQR}^{\ast(k)}_{t,i}$ is the interquartile range of the column over the
  earlier rows whose $k$-interval average holds only observed intervals (the availability mask eroded over about
  $\lceil k \cdot \rhorNBars / \rhdmBarsPerSession \rceil$ rows of the sample). When at least \rhdmModeMinPrints{}
  earlier rows exist and one value accounts for at least \rhdmModeShare\% of them, rows at that value are left out.
  The floor is re-evaluated every \rhdmIqrGridStep{} rows and held in between, and it equals 1 until
  \rhdmIqrMinRows{} such rows exist. It guards against a window that holds mostly imputed values, whose interquartile
  range collapses.
\item \textbf{Columns.} The availability averages $\bar a^{(k)}$ keep their $[0, 1]$ scale. The HAR components and the
  calendar columns are scaled by the same rule without the causal floor ($\mathrm{IQR}^{\ast(k)} = 0$).
\item \textbf{Late start.} For a series whose first print comes after the first day of the sample, an interval is live
  for horizon $k$ when its $k$-interval average holds no interval before the day of the first print. The standardized
  value $z^{(k)}_{t,i}$ is set to 0 on every row whose window $\mathcal W_{t,i}$ holds a row that is not live; since
  liveness only switches on, these are the first $L + W$ rows, $L$ the first live row. A late series therefore enters
  a forecast only once a full window of live averages exists, and because the rows before that point stay at zero,
  the estimation windows of the following years hold few rows on which the series varies (none at the first such
  forecast). The availability averages are not affected.
\item \textbf{Live shares on the option days.} On the closing bars of the \rheSzDays{} post-warmup option days of
  Section~\ref{sec:rh_econ}, the value averages are live on \rhlvVixOptRange\% of bars for the VIX,
  \rhlvVixThreeOptRange\% for VIX3M, \rhlvStAttOptRange\% for StockTwits attention, \rhlvStSentOptRange\% for its
  sentiment and tag count, \rhlvVdOptRange\% for volatility demand and \rhlvVvixOptRange\% for VVIX (ranges over $k$).
  The first prints are in \rhlvVixFirst{} (VIX), \rhlvVixThreeFirst{} (VIX3M), \rhlvStAttFirst{} (StockTwits
  attention), \rhlvStSentFirst{} (sentiment and tag count), \rhlvVdFirst{} (volatility demand) and \rhlvVvixFirst{}
  (VVIX); the one-interval averages turn live in \rhlvVixLiveYear, \rhlvVixThreeLiveYear, \rhlvStAttLiveYear,
  \rhlvStSentLiveYear, \rhlvVdLiveYear{} and \rhlvVvixLiveYear{}, respectively.
\item \textbf{No observation.} After the late-start rule, $z^{(k)}_{t,i} = 0$ on every row whose $k$-interval average
  holds no observed interval ($\bar a^{(k)}_{t,i} = 0$); the scaler's statistics are computed before either rule and
  are not changed by them.
\item \textbf{Why the linear models standardize again.} The step equalizes each column's interquartile range, not its
  variance: the ratio of a column's standard deviation to its interquartile range varies with its tails, its skewness
  and its share of zeros, and each row is scaled with its own trailing window, so the columns of an estimation window
  are not standardized on that window. The step serves the data: it limits the influence of outliers and keeps each
  series comparable over time. The penalized linear models therefore also standardize every column within their own
  estimation window (Appendix~\ref{app:rh_methods_linear}), so that the penalty treats every column alike.
\end{itemize}

% ----------------------------------------------------------------------------------------------------------------------
% From 03_data.tex (last item of the data cleaning) and 04_features.tex (the closing-interval clauses of steps 1, 2, 4
% and 5 and the full-sample list of the causal-construction item), 2026-10-26, with the two target-side differences of
% audit fidelity-5 / design-9 / correctness-6 (results/data_audit_2026-10-15/SUMMARY.md, row 6 and gate a).
% APPENDIX / coordinator: per D3 and D9 the closing-interval models of the paper are the CI-HAR rebuilt under the current
% data chain (results/ci_har_rebuild_2026-10-26/RULE.txt: no exogenous series; B_REGULAR_BARS_ONLY,
% DROP_EARLY_CLOSE_SESSIONS and BURNIN_IN_BAR_SPACE on) and the CLOSE arm (the 16:00 rows of the main design), both on
% the data chain of this appendix. This subsection applies only to forecasts of the earlier closing-interval model set
% (the design cache of 6 October 2026, every data-fix flag off); DROP IT (and its \ref's) if no such forecast remains
% in the paper.
% 2026-10-26 (coordinator, D3/D9): parked -- the earlier closing-interval model set is no longer reported (the CI-HAR is rebuilt on the current data chain; the closing-interval models are the three-way CLOSE arm), so its as-fitted data disclosure goes. Was:
% \subsection{Earlier closing-interval forecasts: data as fitted}
% \label{app:rh_dp_asfitted}
% 2026-10-26 (coordinator: the paper keeps its bullet structure): the lead sentences as bold-lead items, the list nested.
% \begin{itemize}
% \item \textbf{Scope.} The rebuilt closing-interval HAR benchmark and the models trained on the closing half hour
%   (Appendix~\ref{app:rh_closing}) use the data of this appendix.
% \item \textbf{Earlier forecasts, as fitted.} Forecasts of the earlier closing-interval model set, where reported, were
%   fitted before these data fixes:
% \begin{itemize}
% \item no print is set to missing by the domain rule; the implied volatility indices are unobserved after
%   \rhdsVixVendorLast{} and the FOMC series after \rhdsFomcFileLast; the forward fill has no cap; and every unobserved
%   interval is imputed with the median of the series over the full sample;
% \item no average is set to zero for want of observations or before a series starts;
% \item the target's diurnal factor includes the intervals of the early-close sessions after the 13:00~ET close, and the
%   16:00 rows of those sessions are estimation rows;
% \item the standard-deviation floor $\bar\sigma^{(u)}$ of step~1 is the full-sample median of the same standard
%   deviations, and a series is signed when any value in its full sample is negative, which signs the value-weighted
%   quoted spread through a single invalid print;
% \item the inverse hyperbolic sine routing of step~2 uses the share over the full sample, the intervals before the first
%   print included, which routes the same volatility-demand series on every interval; $\kappa_{t,i}$ takes its
%   full-sample median where it is not yet defined; and the FOMC statement-day indicator is logged and winsorized;
% \item the FOMC proximity measures count to and from the statements issued, without the March 2020 rule of
%   Appendix~\ref{app:rh_dp_fomc}, and are winsorized;
% \item the activity indicator is read on the forward-filled series with a full-sample rule: a series is mode-inflated
%   when its full-sample mode accounts for at least \rhdmModeShare\% of all intervals, and
%   $\chi_{t,i} = \mathbf{1}\{u_{t,i} \ne \operatorname{mode}(u)\}$. Under that rule the \rhdmNActive{} series marked in
%   Table~\ref{tab:rh_dm_dictionary} carry the indicator (the number of observations, the \bkSizeVd{} volatility-demand
%   series and the FOMC statement-day indicator) and, in the fitted forecasts, so does the FOMC statement-interval
%   indicator, since dropped from the predictors. Read after the fill, the indicator of a series whose feed has stopped
%   keeps the value of its last print, so the FOMC statement-day indicator equals 1 on every interval after
%   \rhdsFomcFileLast.
% \end{itemize}
% \end{itemize}
